\section{Erd\H{o}s Problem \#43 (Round 8)}

\subsection{1) ROUND-8 OBJECTIVE}

\textbf{Path chosen: (C) --- obstruction/correction for the remaining open part.}

Round~7 ruled out the entire \emph{finite positive-combination} version of the Tao-style kernel method, but it still left one genuine loophole: perhaps one could beat the Round~7 barrier by using a \emph{single general nonnegative positive-definite weight} that is not presented in advance as a nonnegative-generator autocorrelation, or by hiding such a weight inside a more elaborate kernel package.

The most promising Round~8 direction is therefore to close that loophole completely.  I prove a \emph{universal positive-definite-kernel barrier}: among all finitely supported nonnegative positive-definite weights, the unique optimal choice is the discrete Fej\'er kernel, and the best bound obtainable is exactly the same interval-kernel bound already visible in Rounds~6--7.  This strictly strengthens Round~7 by removing both the ``nonnegative generator'' and ``finite package'' restrictions.

\subsection{2) Round-7 FOUNDATION USED}

I explicitly use the following Round~7 results.
\begin{itemize}
\item \textbf{Round~7 pointwise input.} If $A,B\subset[N]$ are Sidon and $(A-A)\cap(B-B)=\{0\}$, then, writing
\[
\Delta_A:=\mathbf 1_A*\mathbf 1_{-A},\qquad \Delta_B:=\mathbf 1_B*\mathbf 1_{-B},
\]
one has
\[
\Delta_A(n)+\Delta_B(n)\le (|A|+|B|)\mathbf 1_{\{0\}}(n)+1
\qquad (n\in\mathbb Z).
\]
\item \textbf{Round~7 finite positive-combination barrier.} Any finite nonnegative linear combination of interval-kernel inequalities is no stronger than the best single interval kernel.
\item \textbf{Round~4 asymptotic lower family.} There exist infinitely many feasible pairs with
\[
|A|+|B|=(\sqrt 2+o(1))\sqrt N.
\]
This will be used to quantify the strength of the new barrier.
\end{itemize}

\subsection{3) NEW INSIGHT / TOOL (ROUND-8)}

There are two genuinely new ingredients in this round.

\paragraph{(i) Fej\'er--Riesz spectral factorisation closes the last loophole in the positive-kernel method.}
I work with an \emph{arbitrary} finitely supported weight $w:\mathbb Z\to [0,\infty)$ whose Fourier transform
\[
\widehat w(\theta):=\sum_{n\in\mathbb Z} w(n)e^{2\pi i n\theta}
\]
is nonnegative for all $\theta\in\mathbb R/\mathbb Z$ (equivalently, $w$ is positive-definite).  By the classical Fej\'er--Riesz theorem, such a $w$ factors as
\[
w=\psi*\widetilde\psi
\]
for some finitely supported complex $\psi$.  This lets me re-run the Round~6/7 argument for the full class of nonnegative positive-definite weights, not just the special weights built from nonnegative generators.

\paragraph{(ii) Exact optimisation of the whole positive-definite-weight framework.}
If $w$ is supported in $\{-(L-1),\dots,L-1\}$, the best upper bound obtainable from $w$ is
\[
B_w(N,T):=(N+L-1)\left(1+T\frac{w(0)}{\sum_n w(n)}\right),
\qquad T:=|A|+|B|.
\]
I prove
\[
\sum_n w(n)\le L\,w(0),
\]
with equality if and only if $w$ is a positive scalar multiple of the discrete Fej\'er kernel
\[
F_L(n):=\max\{L-|n|,0\}.
\]
Hence the exact optimum over \emph{all} admissible positive-definite weights is
\[
\mathcal B_{\mathrm{pd}}(N,T)=\min_{L\ge 1}(N+L-1)\left(1+\frac{T}{L}\right),
\]
and it is uniquely attained (up to scale) by the Fej\'er kernels.  This strictly strengthens Round~7.

\subsection{4) ATTACK PLAN (ROUND-8)}

\textbf{Exact gap left after Round~7.}
Round~7 exhausted finite positive combinations of the previous interval/autocorrelation kernels, but not the larger class of all nonnegative positive-definite weights.  In principle, such a weight could arise from a sign-changing or complex generator after spectral factorisation, so Round~7 did not yet fully rule out the broader Tao-style framework.

\textbf{Claims proved in this round.}
\begin{enumerate}
\item Every admissible positive-definite weight $w$ yields an upper bound
\[
S:=|A|^2+|B|^2\le B_w(N,T):=(N+L-1)\left(1+T\frac{w(0)}{\sum_n w(n)}\right),
\]
where $L$ is the support-length parameter of $w$.
\item For every such $w$, one has
\[
\sum_n w(n)\le L\,w(0),
\]
with equality exactly for positive multiples of the Fej\'er kernel.
\item Consequently,
\[
\inf_w B_w(N,T)=\min_{L\ge 1}(N+L-1)\left(1+\frac{T}{L}\right),
\]
so the whole positive-definite-weight method is \emph{exactly} optimised by the same Fej\'er/interval kernel already present in Round~6.
\end{enumerate}

This overcomes the final loophole in the positive-kernel approach and rules out an even broader class of proof strategies than Round~7.

\subsection{5) WORK (ROUND-8)}

\subsubsection{5.1 Admissible weights and Fej\'er--Riesz factorisation}

Let $A,B\subset[N]$ be Sidon sets with $(A-A)\cap(B-B)=\{0\}$, and write
\[
a:=|A|,\qquad b:=|B|,\qquad S:=a^2+b^2,\qquad T:=a+b.
\]

Call a function $w:\mathbb Z\to[0,\infty)$ \emph{admissible of span $L$} if:
\begin{itemize}
\item $w$ is supported in $\{-(L-1),\dots,L-1\}$,
\item $\widehat w(\theta)\ge 0$ for all $\theta\in\mathbb R/\mathbb Z$.
\end{itemize}
The second condition is equivalent to positive-definiteness of $w$.

We will use the following classical theorem.

\begin{theorem}[Fej\'er--Riesz]
\label{thm:FR}
Let
\[
p(\theta)=\sum_{n=-(L-1)}^{L-1} c_n e^{2\pi i n\theta}
\]
be a trigonometric polynomial that is nonnegative for all $\theta\in\mathbb R/\mathbb Z$.  Then there exists a polynomial
\[
P(z)=\sum_{j=0}^{L-1}\psi(j)z^j
\]
of degree at most $L-1$ such that
\[
p(\theta)=\bigl|P(e^{2\pi i\theta})\bigr|^2.
\]
Equivalently, the coefficients satisfy
\[
c_n=\sum_{j\in\mathbb Z}\psi(j+n)\overline{\psi(j)}.
\]
\end{theorem}

Applied to $p=\widehat w$, Theorem~\ref{thm:FR} gives the following factorisation.

\begin{corollary}
\label{cor:factor}
If $w$ is admissible of span $L$, then there exists a complex-valued function $\psi$ supported in $\{0,1,\dots,L-1\}$ such that
\[
w=\psi*\widetilde\psi,
\qquad \widetilde\psi(n):=\overline{\psi(-n)}.
\]
\end{corollary}

\begin{proof}
Write
\[
\widehat w(\theta)=\sum_{n=-(L-1)}^{L-1} w(n)e^{2\pi i n\theta}.
\]
Since $w$ is admissible, $\widehat w(\theta)\ge 0$ for all $\theta$.  Applying Theorem~\ref{thm:FR} gives
\[
\widehat w(\theta)=\left|\sum_{j=0}^{L-1}\psi(j)e^{2\pi i j\theta}\right|^2.
\]
Comparing Fourier coefficients yields $w=\psi*\widetilde\psi$.
\end{proof}

\subsubsection{5.2 The universal positive-definite-kernel inequality}

\begin{proposition}
\label{prop:pd-bound}
Let $w$ be an admissible weight of span $L$, and set
\[
W:=\sum_{n\in\mathbb Z} w(n).
\]
Then
\begin{equation}
\label{eq:pd-master}
S\le (N+L-1)\left(1+T\frac{w(0)}{W}\right).
\end{equation}
\end{proposition}

\begin{proof}
By Corollary~\ref{cor:factor}, write $w=\psi*\widetilde\psi$ with $\operatorname{supp}(\psi)\subseteq\{0,\dots,L-1\}$.

First, using the Round~7 pointwise estimate,
\[
\Delta_A(n)+\Delta_B(n)\le T\mathbf 1_{\{0\}}(n)+1,
\]
and the fact that $w(n)\ge0$, we obtain
\begin{align}
\sum_n w(n)\Delta_A(n)+\sum_n w(n)\Delta_B(n)
&\le T w(0)+\sum_n w(n)\notag\\
&=T w(0)+W.
\label{eq:upperpd}
\end{align}

On the other hand,
\[
\sum_n w(n)\Delta_A(n)=\sum_n (\psi*\widetilde\psi)(n)(\mathbf 1_A*\mathbf 1_{-A})(n)=\|\mathbf 1_A*\psi\|_2^2,
\]
and similarly
\[
\sum_n w(n)\Delta_B(n)=\|\mathbf 1_B*\psi\|_2^2.
\]
Now $\mathbf 1_A*\psi$ is supported in an interval of length at most $N+L-1$, so Cauchy--Schwarz gives
\[
\|\mathbf 1_A*\psi\|_2^2\ge \frac{\left|\sum_x (\mathbf 1_A*\psi)(x)\right|^2}{N+L-1}
=\frac{a^2\left|\sum_j \psi(j)\right|^2}{N+L-1}.
\]
Likewise,
\[
\|\mathbf 1_B*\psi\|_2^2\ge \frac{b^2\left|\sum_j \psi(j)\right|^2}{N+L-1}.
\]
Adding and using
\[
W=\sum_n w(n)=\left|\sum_j \psi(j)\right|^2,
\]
we obtain
\begin{equation}
\label{eq:lowerpd}
\sum_n w(n)\Delta_A(n)+\sum_n w(n)\Delta_B(n)
\ge \frac{SW}{N+L-1}.
\end{equation}
Combining \eqref{eq:upperpd} and \eqref{eq:lowerpd} yields \eqref{eq:pd-master}.
\end{proof}

\subsubsection{5.3 Fej\'er's inequality and the unique optimiser}

\begin{lemma}[Fej\'er inequality in coefficient form]
\label{lem:fejer-ineq}
If $w$ is admissible of span $L$, then
\begin{equation}
\label{eq:fejer-coeff}
W\le L\,w(0).
\end{equation}
Moreover, equality holds if and only if
\[
w(n)=c\,F_L(n)
\qquad (n\in\mathbb Z)
\]
for some $c>0$, where
\[
F_L(n):=\max\{L-|n|,0\}
\]
is the discrete Fej\'er kernel.
\end{lemma}

\begin{proof}
By Corollary~\ref{cor:factor}, write $w=\psi*\widetilde\psi$ with $\operatorname{supp}(\psi)\subseteq\{0,\dots,L-1\}$.
Then
\[
W=\left|\sum_{j=0}^{L-1}\psi(j)\right|^2,
\qquad
w(0)=\sum_{j=0}^{L-1}|\psi(j)|^2.
\]
Applying Cauchy--Schwarz to the length-$L$ vector $(\psi(0),\dots,\psi(L-1))$ gives
\[
\left|\sum_{j=0}^{L-1}\psi(j)\right|^2\le L\sum_{j=0}^{L-1}|\psi(j)|^2,
\]
which is exactly \eqref{eq:fejer-coeff}.

Equality in Cauchy--Schwarz holds if and only if all $\psi(j)$ have the same phase and the same modulus.  Thus, after multiplying by a unimodular constant, one has
\[
\psi(j)=c^{1/2}
\qquad (0\le j\le L-1)
\]
for some $c>0$.  In that case,
\[
w(n)=\sum_j \psi(j+n)\overline{\psi(j)}=c\,\#\{j:0\le j\le L-1,\ 0\le j+n\le L-1\}=c\,F_L(n).
\]
Conversely, $w=cF_L$ is generated by the constant vector $\psi\equiv c^{1/2}$ on $\{0,\dots,L-1\}$, so equality holds.
\end{proof}

\begin{theorem}[Universal positive-definite-kernel barrier]
\label{thm:universal-barrier}
Let $\mathcal B_{\mathrm{pd}}(N,T)$ denote the strongest upper bound on $S=|A|^2+|B|^2$ obtainable from the Round~7 pointwise estimate by summing against an arbitrary admissible positive-definite weight and using only the Cauchy--Schwarz lower bound above.  Then
\begin{equation}
\label{eq:Bpd}
\mathcal B_{\mathrm{pd}}(N,T)=\min_{L\ge 1}(N+L-1)\left(1+\frac{T}{L}\right).
\end{equation}
For each fixed $L$, the unique optimiser (up to positive scalar multiple) is the discrete Fej\'er kernel $F_L$.
\end{theorem}

\begin{proof}
Fix an admissible weight $w$ of span $L$.  Proposition~\ref{prop:pd-bound} gives
\[
S\le (N+L-1)\left(1+T\frac{w(0)}{W}\right).
\]
By Lemma~\ref{lem:fejer-ineq},
\[
\frac{w(0)}{W}\ge \frac1L,
\]
so every admissible weight of span $L$ yields a bound at least
\[
(N+L-1)\left(1+\frac{T}{L}\right).
\]
Conversely, taking $w=F_L=\mathbf 1_{[0,L-1]}*\mathbf 1_{-[0,L-1]}$ gives
\[
W=L^2,
\qquad
w(0)=L,
\]
so Proposition~\ref{prop:pd-bound} yields exactly
\[
S\le (N+L-1)\left(1+\frac{T}{L}\right).
\]
Minimising over $L\ge 1$ proves \eqref{eq:Bpd}.  The uniqueness statement for fixed $L$ follows from the equality case in Lemma~\ref{lem:fejer-ineq}.
\end{proof}

\subsubsection{5.4 Consequences and relation to Round~7}

\begin{corollary}
\label{cor:subsumes-r7}
Round~7's finite positive-combination barrier is a special case of Theorem~\ref{thm:universal-barrier}.
\end{corollary}

\begin{proof}
A finite nonnegative combination of Round~7 kernels has the form
\[
w=\sum_{i=1}^m \lambda_i (\phi_i*\widetilde\phi_i),
\qquad \lambda_i\ge0.
\]
Such a $w$ is entrywise nonnegative, finitely supported, and positive-definite because
\[
\widehat w(\theta)=\sum_{i=1}^m \lambda_i\,|\widehat{\phi_i}(\theta)|^2\ge 0.
\]
Therefore $w$ is admissible, so Theorem~\ref{thm:universal-barrier} applies directly.
\end{proof}

\begin{corollary}[The barrier is genuinely of order $N^{3/4}$ on near-extremal families]
\label{cor:n34}
For the infinite family from Round~4 with
\[
T=|A|+|B|=(\sqrt2+o(1))\sqrt N,
\]
the optimum of the whole positive-definite-weight method satisfies
\[
\mathcal B_{\mathrm{pd}}(N,T)
\ge N+T-1+2\sqrt{T(N-1)}
= N+2^{5/4}N^{3/4}+o(N^{3/4}).
\]
In particular, no proof inside this framework can yield a uniform $N+O(1)$ bound for $S$, and hence none can prove the original $O(1)$-gap conjecture.
\end{corollary}

\begin{proof}
From Theorem~\ref{thm:universal-barrier},
\[
\mathcal B_{\mathrm{pd}}(N,T)=\min_{L\ge 1}\left(N+T+L-1+\frac{T(N-1)}{L}\right).
\]
By AM--GM,
\[
L+\frac{T(N-1)}{L}\ge 2\sqrt{T(N-1)}.
\]
Hence
\[
\mathcal B_{\mathrm{pd}}(N,T)\ge N+T-1+2\sqrt{T(N-1)}.
\]
Substituting $T=(\sqrt2+o(1))\sqrt N$ gives
\[
2\sqrt{T(N-1)}=(2^{5/4}+o(1))N^{3/4},
\]
and the result follows.
\end{proof}

\subsection{6) ADVERSARIAL VERIFICATION}

\paragraph{(i) Did Round~8 really strengthen Round~7?}
Yes.  Round~7 treated finite positive combinations of pre-chosen autocorrelation kernels generated by nonnegative functions.  Round~8 removes both restrictions: the weight may now be any finitely supported nonnegative positive-definite sequence, equivalently any nonnegative trigonometric polynomial in the frequency side.  Theorem~\ref{thm:universal-barrier} therefore subsumes Round~7.

\paragraph{(ii) Hidden hypothesis in Fej\'er--Riesz.}
Theorem~\ref{thm:FR} applies because $\widehat w$ is a one-variable nonnegative trigonometric polynomial of degree at most $L-1$.  This exactly matches the admissible-weight setup.

\paragraph{(iii) Why is $W=\sum_n w(n)$ nonzero?}
If $w\not\equiv 0$ and $w(n)\ge0$ for all $n$, then $W>0$.  The zero weight yields no information and is excluded from optimisation.

\paragraph{(iv) Equality case in Lemma~\ref{lem:fejer-ineq}.}
The equality statement uses only the equality case of Cauchy--Schwarz.  Once all coefficients $\psi(0),\dots,\psi(L-1)$ have the same phase and modulus, the resulting autocorrelation is exactly a positive multiple of $F_L$.  There is no additional hidden symmetry assumption.

\paragraph{(v) Does this settle the first question?}
No.  It settles only a \emph{methodological} question: every proof that uses only the pointwise bound and a single nonnegative positive-definite weight (equivalently any finite positive combination of such weights) is now exhausted, and its optimum is exactly known.  The first question itself remains open.

\subsection{7) FINAL}

\textbf{UNRESOLVED (BUT STRICTLY ADVANCED).}

Round~8 proves a genuinely stronger barrier theorem than Round~7.  The whole nonnegative positive-definite-kernel method is now optimised exactly:
\[
\mathcal B_{\mathrm{pd}}(N,T)=\min_{L\ge 1}(N+L-1)\left(1+\frac{T}{L}\right),
\]
and the unique optimisers are the discrete Fej\'er kernels.  Thus the interval/Fej\'er kernel was not merely a convenient choice in earlier rounds --- it is the best possible kernel in the entire positive-definite-weight framework.  This rules out an even broader proof strategy for the first question.

\subsection{8) COMPLETION ESTIMATE (MANDATORY)}

\textbf{COMPLETION: 93\%}

\subsection{9) REFERENCES}

\begin{itemize}
\item Round~7 write-up supplied in this task.
\item M. A. Dritschel and J. Rovnyak, \emph{The Operator Fej\'er--Riesz Theorem}, for the classical one-variable Fej\'er--Riesz factorisation statement used here.
\item T. F. Bloom, \emph{Erd\H{o}s Problem \#43}, problem page and discussion thread.
\end{itemize}
